Put \(N=N_m\), \(K=K_{\epsilon}\),
\[
 h(q)=(g_m(q),f(q)),\qquad S=h(K),\qquad H=\operatorname{conv}(S).
\tag{1}
\]
The set \(K\) is the intersection of the compact simplex with the closed inequalities \(\epsilon\leq p(q)\leq1-\epsilon\), so it is compact.  Both denominators in \(f\) are at least \(\epsilon>0\) on \(K\); hence \(f\), \(g_m\), and \(h\) are continuous.  The multinomial theorem gives, for every \(q\in K\),
\[
 \mathbf 1^{\top}g_m(q)=\sum_{c\in\mathcal I_m}B_c^{(m)}(q)
 =\left(\sum_{a,y}q_{ay}\right)^m=1.
\tag{2}
\]

We first establish the dimensions needed for endpoint compression.  The set
\[
 K^{\circ}:=\{q\in\Delta_3:q_{ay}>0\text{ for every }(a,y),\ 
 \epsilon<p(q)<1-\epsilon\}
\tag{3}
\]
is a nonempty relatively open subset of the simplex: the uniform four-cell vector belongs to it because \(0<\epsilon<1/2\).  Suppose
\[
 P(q):=\sum_{c\in\mathcal I_m}a_cB_c^{(m)}(q)=0\qquad(q\in K).
\tag{4}
\]
The left side is a homogeneous polynomial of degree \(m\).  Substituting \(q_{11}=1-q_{00}-q_{01}-q_{10}\) makes it a polynomial in three variables that vanishes on the nonempty open set corresponding to \(K^{\circ}\); successive use of the elementary fact that a univariate polynomial with an interval of roots is zero shows that this three-variable polynomial is identically zero.  Thus \(P=0\) on the entire affine hyperplane \(\sum q_{ay}=1\).  If \(v\in\mathbb R^4\) and \(s:=\sum v_{ay}\neq0\), homogeneity yields
\[
 P(v)=s^mP(v/s)=0.
\tag{5}
\]
The set \(\{v:s\neq0\}\) is open, so the polynomial \(P\) is identically zero on \(\mathbb R^4\).  The distinct degree-\(m\) monomials are linearly independent and every Bernstein normalization is nonzero, whence all \(a_c=0\).  Therefore the \(N\) Bernstein coordinates are linearly independent as functions on \(K\).

Equation (2) places \(g_m(K)\) in the hyperplane
\[
 A_0:=\{x\in\mathbb R^N:\mathbf 1^{\top}x=1\}.
\tag{6}
\]
If \(\operatorname{aff}g_m(K)\) were a proper affine subspace of \(A_0\), some affine functional nonconstant on \(A_0\) would be constant on \(g_m(K)\).  After using (2) to absorb its constant term, this would give a nonzero linear combination of the Bernstein coordinates that vanishes on \(K\), contradicting (4)--(5).  Hence
\[
 \operatorname{aff}g_m(K)=A_0,\qquad
 \dim\operatorname{aff}g_m(K)=N-1.
\tag{7}
\]
The set \(H\) lies in \(A_0\times\mathbb R\), of dimension \(N\), and its projection contains \(g_m(K)\); hence (7) gives \(\dim\operatorname{aff}H\geq N-1\).  If \(\dim\operatorname{aff}H=N-1\), projection onto the first \(N\) coordinates would map the translation space of \(\operatorname{aff}H\) bijectively onto that of \(A_0\), by (7).  The last coordinate on \(\operatorname{aff}H\) would consequently be an affine function of the first coordinates.  Thus there would be \(\lambda\in\mathbb R^N\) such that
\[
 f(q)=\lambda^{\top}g_m(q)\qquad(q\in K),
\tag{8}
\]
where an affine constant is again absorbed using (2).  Equation (8) says \(f\in V_m\), contrary to the finite-plate nonpolynomial lemma.  It follows that
\[
 \dim\operatorname{aff}H=N,\qquad
 \operatorname{aff}H=A_0\times\mathbb R.
\tag{9}
\]

We next prove compactness and the attainable-pair identity.  Every finite convex representation in \(H\) with more than \(N+1\) positive weights is affinely dependent because of (2).  Thus, if
\[
 z=\sum_{j=1}^k w_js_j,qquad w_j>0,qquad k>N+1,
\]
there are coefficients \(a_j\), not all zero, satisfying \(\sum_ja_j=0\) and \(\sum_ja_js_j=0\).  Both signs occur.  With
\[
 \rho=\min_{j:a_j>0}\frac{w_j}{a_j},\qquad w'_j=w_j-\rho a_j,
\tag{10}
\]
the weights \(w'_j\) are nonnegative, sum to one, represent the same \(z\), and at least one is zero.  Iteration leaves at most \(N+1\) positive weights.  Hence \(H\) is the image of the compact set \(\Delta_N\times K^{N+1}\) under the continuous barycenter map, and is compact.

Every point of \(H\) is the pair of integrals under its finite atomic representing measure.  Conversely, for \(\nu\in\mathcal P(K)\), put \(z=\int h\,d\nu\).  If \(z\notin H\), compactness supplies a closest \(z_0\in H\).  For every \(x\in H\), differentiation at zero along the segment \(z_0+u(x-z_0)\) gives
\[
 (z-z_0)^{\top}(x-z_0)\leq0.
\tag{11}
\]
Taking \(x=h(q)\), integrating (11), and using \(z=\int h\,d\nu\) yields \(\lVert z-z_0\rVert^2\leq0\), a contradiction.  Therefore
\[
 H=
 \left\{\left(\int g_m\,d\nu,\int f\,d\nu\right):
 \nu\in\mathcal P(K)\right\},
\tag{12}
\]
and every pair in (12) has a representing measure with at most \(N+1\) atoms.

Fix \(b\in\mathcal B_{m,\epsilon}\).  By the definition of the moment body and (12),
\[
 H_b:=\{t:(b,t)\in H\}
 =\left\{\int f\,d\nu:\nu\in\mathcal F_{m,\epsilon}(b)\right\}.
\tag{13}
\]
This section is nonempty, compact, and convex; hence it is a closed interval.  By the definitions of \(L_m\), \(U_m\), and \(\Theta_{m,\epsilon}\),
\[
 H_b=[L_m(b),U_m(b)]=\Theta_{m,\epsilon}(b),
\tag{14}
\]
and both endpoints are attained.

We now improve the support bound at every endpoint.  Let \(z_-=(b,L_m(b))\).  If \(z_-\) were in \(\operatorname{relint}H\), then (9) would put \(z_--(0,\delta)\) in \(H\) for all sufficiently small \(\delta>0\), contradicting (14).  Since \(H\) is closed in its affine hull, \(z_-\in H\setminus\operatorname{relint}H=\operatorname{relbd}H\).  Choose \(z_k\in\operatorname{aff}H\setminus H\) with \(z_k\to z_-\), let \(h_k\) be a closest point of the compact set \(H\) to \(z_k\), and set
\[
 n_k=\frac{h_k-z_k}{\lVert h_k-z_k\rVert}.
\]
Because \(\operatorname{dist}(z_k,H)\leq\lVert z_k-z_-\rVert\), one has \(h_k\to z_-\).  The nearest-point inequality along every segment from \(h_k\) to \(y\in H\) is
\[
 n_k^{\top}(y-h_k)\geq0.
\tag{15}
\]
Along a subsequence, \(n_k\) converges to a unit vector \(n\) in the translation space of \(\operatorname{aff}H\), and (15) gives
\[
 n^{\top}(y-z_-)\geq0\qquad(y\in H).
\tag{16}
\]
Thus
\[
 F_-:=\{y\in H:n^{\top}(y-z_-)=0\}
\tag{17}
\]
is an exposed face containing \(z_-\).  It is proper, since otherwise \(\operatorname{aff}H\) would lie in the hyperplane normal to the nonzero vector \(n\), contrary to (9).  Therefore \(\dim\operatorname{aff}F_-\leq N-1\).

Represent \(z_-\) by at most \(N+1\) points \(s_j\in S\), as established above.  Every term \(n^{\top}(s_j-z_-)\) is nonnegative by (16), while their positive-weighted sum is zero.  Hence every support point lies in \(S\cap F_-\).  Applying the weight-elimination argument (10) inside the affine space \(\operatorname{aff}F_-\) reduces the representation to at most
\[
 \dim\operatorname{aff}(S\cap F_-)+1\leq N
\tag{18}
\]
points.  Replacing \(z_-\) by \(z_+=(b,U_m(b))\) and using an exterior sequence above the section gives the same conclusion for the upper endpoint.  By the converse part of the SCM moment-reduction lemma, each endpoint measure yields a compatible member of \(\mathfrak M_{m,\epsilon}\) with the observed count law \(b\), causal ATE equal to the endpoint, and at most \(N=N_m\) latent laws.  Notice that (18) concerns vertical-section endpoints; the \(N+1\) cap for an arbitrary point of the full hull in (12) is unchanged.

It remains to prove the asserted attained duality for \(b\in\operatorname{relint}(\mathcal B_{m,\epsilon})\).  Put \(A=\operatorname{aff}(\mathcal B_{m,\epsilon})\), let \(V\) be its translation space, and define in \(A\times\mathbb R\)
\[
 E_-:=H+\{(0,s):s\geq0\}
 =\{(x,t):x\in\mathcal B_{m,\epsilon},\ t\geq L_m(x)\}.
\tag{19}
\]
The set is closed: in any convergent sequence \(h_k+(0,s_k)\), compactness of \(H\) supplies a convergent subsequence of \(h_k\), after which \(s_k\) converges to a nonnegative limit.  Project \((b,L_m(b)-k^{-1})\) onto \(E_-\) and pass to a limiting normalized nearest-point vector as in (15).  This gives \((a,\beta)\in V\times\mathbb R\), of unit norm, such that
\[
 a^{\top}(x-b)+\beta\{t-L_m(b)\}\geq0
 \qquad((x,t)\in E_-).
\tag{20}
\]
Upward closure forces \(\beta\geq0\).  If \(\beta=0\), a relative ball about \(b\) in \(A\) and the choices \(x=b\pm\delta v\), for arbitrary \(v\in V\), imply \(a^{\top}v=0\).  Hence \(a=0\), contradicting unit norm.  Thus \(\beta>0\), and we divide (20) by \(\beta\).

Define
\[
 \lambda^-=-a+\{L_m(b)+a^{\top}b\}\mathbf 1.
\tag{21}
\]
Equations (2) and (20), evaluated at \(h(q)\in H\), give
\[
 (\lambda^-)^{\top}g_m(q)
 =L_m(b)-a^{\top}\{g_m(q)-b\}\leq f(q),
 \qquad
 (\lambda^-)^{\top}b=L_m(b).
\tag{22}
\]
For every feasible minorant \(\eta\) and every \(\nu\in\mathcal F_{m,\epsilon}(b)\),
\[
 \eta^{\top}b=\int\eta^{\top}g_m\,d\nu
 \leq\int f\,d\nu.
\tag{23}
\]
Taking the infimum in (23) proves weak duality, and (22) proves equality and attainment for the lower program.

Similarly, for
\[
 E_+:=H-\{(0,s):s\geq0\}
 =\{(x,t):x\in\mathcal B_{m,\epsilon},\ t\leq U_m(x)\},
\]
projection from \((b,U_m(b)+k^{-1})\) gives, after scaling, an inequality
\[
 a_+^{\top}(x-b)-\{t-U_m(b)\}\geq0\qquad((x,t)\in E_+).
\tag{24}
\]
Consequently
\[
 \lambda^+=a_++\{U_m(b)-a_+^{\top}b\}\mathbf 1
\tag{25}
\]
satisfies
\[
 (\lambda^+)^{\top}g_m(q)\geq f(q),
 \qquad (\lambda^+)^{\top}b=U_m(b).
\tag{26}
\]
Integration of any feasible majorant gives the reverse weak-duality inequality.  Equations (22)--(26) therefore prove
\[
 D_m(b)=(L_m(b),U_m(b))
\tag{27}
\]
with both dual extrema attained.

Finally, for the lower optimizer define
\[
 T_-:=\{q\in K:(\lambda^-)^{\top}g_m(q)=f(q)\}.
\]
If \(\nu_-\) attains the lower endpoint, (22) gives
\[
 0=\int_K\{f(q)-(\lambda^-)^{\top}g_m(q)\}\,d\nu_-(q).
\tag{28}
\]
The integrand is continuous and nonnegative, so \(\nu_-(T_-)=1\).  Hence \(b\in\operatorname{conv}g_m(T_-)\).  If \(r=\dim\operatorname{aff}g_m(T_-)\), the elimination in (10) gives an endpoint measure on at most \(r+1\) contact points, because contact makes its objective equal to \((\lambda^-)^{\top}b=L_m(b)\).  The identical argument for the upper contact set proves the upper assertion.  This completes all primal, causal-witness, dual, and contact-support claims.